(1) 理想磁流体的应力张量为 $\mathbf T=-(p+B^2/(2\mu_0))\mathbf I+\mathbf B\mathbf B/\mu_0$。题设磁场完全切于界面，法向张力部分为零，所以没有额外表面力时两侧总压力满足 $p_1+B_1^2/(2\mu_0)=p_2+B_2^2/(2\mu_0)$。
(2) 整理得 $p_1-p_2=(B_2^2-B_1^2)/(2\mu_0)$。若 $B_1>B_2$，第一侧的磁压力更大，必须由相应较小气体压强补偿，即 $p_1<p_2$；不能把磁压力和气体压力在平衡时都简单看成同向增大。
(3) 让单位面积界面向区域2移动微小长度 $dx$，总应力的可用机械功差为 $dW=[p_1+B_1^2/(2\mu_0)-p_2-B_2^2/(2\mu_0)]dx$。利用界面压力条件，该差严格等于零，符合静态平衡不发生无驱动自发位移的要求。
(4) 若两侧磁场强度相同，则 $p_1=p_2$；若全无磁场，也回到普通流体界面压力连续。若存在非零法向磁场，张量项 $\mathbf B\mathbf B/\mu_0$ 贡献法向拉应力，必须保留它，不能只用标量磁压简单相加。
